\subsection{无限整环上的多项式函数}

假设 A 是一个整环。设 I 是一个集合， \((\mathrm{H}_{i})_{i\in I}\) 为 A 的无限子集的一个族，且 \(H=\prod_{i\in I}H_{i}\subset A^{I}\) 。若 f 是非零

的元素 \(\mathbf{A}[(\mathbf{X}_i)_{i\in I}]\) ，并设 \(\mathrm{H}_f\) 是所有 \(\mathbf{x}\in \mathrm{H}\) 中满足 \(f(\mathbf{x})\neq 0\) 者组成的集合，则 \(\mathrm{H}\) 与 \(\mathrm{H}_f\) 等势。

\begin{enumerate}
\def\labelenumi{\alph{enumi})}
\tightlist
\item
  首先假设 I 是有限的，并令 \(n = \text{Card I}\) 。该命题对 \(n = 0\) 是显然的；我们将对 \(n\) 作归纳来证明它。选取 I 的一个元素 \(i_0\) ，并令 \(J = I - \{i_0\}\) , \(B = A[(X_i)_{i \in J}]\) 。由于 \(f \neq 0\) ，我们可以写出 \(f = \sum_{k=0}^{m} g_k X_{i_0}^k\) ，其中 \(g_0, \ldots, g_m \in \mathbf{B}\) 且 \(g_m \neq 0\) 。由归纳假设，集合 \(\mathbf{K}\) （由所有 \(\mathbf{x} \in \prod_{i \in J} \mathrm{H}_i\) 中满足 \(g_m(\mathbf{x}) \neq 0\) 者组成）与 \(\prod_{i \in J} \mathrm{H}_i\) 等势。对 \(\mathbf{x} \in \mathbf{K}\) ，多项式
\end{enumerate}

\[
h \left(\mathrm{X} _ {i _ {0}}\right) = \sum_ {k = 0} ^ {m} g _ {k} (\mathbf {x}) \mathrm{X} _ {i _ {0}} ^ {k} \in \mathrm{A} \left[ \mathrm{X} _ {i _ {0}} \right]
\]

非零。根据定理 1（IV，第 16 页），所有 \(\alpha \in \mathrm{H}_{i_0}\) 中满足 \(h(\alpha) \neq 0\) 者组成的集合与 \(\mathrm{H}_{i_0}\) 等势，由此

\[
\operatorname{Card} \mathrm{H} \geqslant \operatorname{Card} \mathrm{H} _ {f} \geqslant (\operatorname{Card} \mathrm{K}). (\operatorname{Card} \mathrm{H} _ {i _ {0}}) = \operatorname{Card} \mathrm{H},
\]

因此 \(\mathbf{H} = \mathbf{Card}\mathbf{H}_f\)

\begin{enumerate}
\def\labelenumi{\alph{enumi})}
\setcounter{enumi}{1}
\tightlist
\item
  在一般情形，存在有限子集 \(\mathbf{I}'\) （ \(\mathbf{I}\) 的）使得 \(f \in \mathbf{A}[(\mathbf{X}_i)_{i \in \mathbf{I}'}]\) 。设 \(\mathbf{H}_f'\) 是所有 \(\mathbf{x} \in \prod_{i \in \mathbf{I}'} \mathbf{H}_i\) 中满足 \(f(\mathbf{x}) \neq 0\) 者组成的集合，则 \(\mathbf{H}_f = \mathbf{H}_f' \times \left( \prod_{i \in \mathbf{I} - \mathbf{I}'} \mathbf{H}_i \right)\) ，只需将证明的第一部分应用于 \(\mathbf{H}_f'\) 。
\end{enumerate}

推论 1. --- 我们保留命题 9 的假设和记号。若 I 非空，则 \(H_{f}\) 是无限的。

推论 2. --- 设 A 是无限整环，或 A 是无限域上的代数。对每一个 \(f \in \mathbf{A}[(\mathbf{X}_i)_{i \in I}]\) ，设 \(\tilde{f}: \mathbf{A}^{\mathrm{I}} \to \mathbf{A}\) 为由 \(f\) 定义的多项式函数（IV，第 4 页）。于是映射 \(f \mapsto \tilde{f}\) 是单射的。

当 \(\mathbf{A}\) 是一个无限整环，则该推论立即由命题 9 得出。假设 \(\mathbf{A}\) 是某个无限域上的代数，该域记作 \(k\) 。设 \(f = \sum_{\nu \in \mathbb{N}^{(I)}} \alpha_{\nu} X^{\nu}\)

是 \(\mathbf{A}[(\mathbf{X}_i)_{i\in I}]\) 中的非零元素；则存在 \(\nu_0\in \mathbf{N}^{(I)}\) 使得 \(\alpha_{\nu_0}\neq 0\) ，以及一个 \(k\) -线性形式 \(\varphi\) （ \(\mathbf{A}\) 上的）使得 \(\varphi (\alpha_{\nu_0})\neq 0\) 。设 \(g = \sum_{\nu \in \mathbf{N}^{(I)}}\varphi (\alpha_{\nu})\mathbf{X}^{\nu}\in k[(\mathbf{X}_i)_{i\in I}]\) ；我们有 \(g\neq 0\) ，因此存在 \(\mathbf{x}\in k^{\mathrm{I}}\) 使得

\[
g (\mathbf {x}) \neq 0
\]

\[
\varphi (f (\mathbf {x})) = g (\mathbf {x}) \neq 0
\]

\[
f (\mathbf {x}) \neq 0
\]

当 A 是无限整环，或当 A 是无限域上的代数时，我们通常将 f 与 \(\tilde{f}\) 等同。

假设 A 是有限的，并令 \(f(X) = \prod_{\alpha \in A} (X - \alpha)\) ，则 \(f \neq 0\) ，但 \(\tilde{f} = 0\) 。其他例子见 IV，第 88 页，习题 7 和 8。

定理 2（代数恒等式延拓原理）。--- 假设 A 是无限整环。设 \(g_1, \ldots, g_m\) , \(f\) 为 \(\mathbf{A}[(\mathbf{X}_i)_{i \in I}]\) 的元素，并假设以下假设成立：

\begin{enumerate}
\def\labelenumi{\alph{enumi})}
\item
  \(g_{1} \neq 0, \ldots, g_{m} \neq 0\) ;
\item
  对所有 \(\mathbf{x} \in \mathbf{A}^{I}\) （满足 \(g_1(\mathbf{x}) \neq 0, \ldots, g_m(\mathbf{x}) \neq 0\) ），有 \(f(\mathbf{x}) = 0\) 。则 \(f = 0\) 。
\end{enumerate}

因为如果 \(f \neq 0\) ，我们有 \(fg_1 \ldots g_m \neq 0\) （IV，第 9 页，命题 8），因此存在 \(\mathbf{x} \in \mathbf{A}^{I}\) 使得 \(f(\mathbf{x})g_1(\mathbf{x}) \ldots g_m(\mathbf{x}) \neq 0\) （IV，第 18 页，推论 2），这与假设矛盾。

附释。--- 设 A 是整环且 \(f \in \mathbf{A}[(\mathbf{X}_i)_{i \in I}]\) 。定理 2 提供了证明 f = 0 的一个便捷方法。只需考虑一个以 A 为子环的无限整环 E；如果我们能证明 \(f((x_i)) = 0\) 对所有 \((x_i) \in \mathrm{E}^{\mathrm{I}}\) 成立（甚至仅对 \((x_i) \in \mathrm{E}^{\mathrm{I}}\) 中那些使有限个给定的非零多项式不消失的点成立），那么就可以推出 f = 0。如果 A 本身不是无限的，例如可以取 E 为环 A{[}X{]} 或其分式域。

一旦我们证明了关系式 \(f = 0\) ，我们显然可以推断出 \(f((y_i)) = 0\) （对所有 \((y_i) \in \mathbf{F}^{I}\) ），其中 \(\mathbf{F}\) 是任意一个含单位元的结合且交换的 \(\mathbf{A}\) -代数；特别地， \(\mathbf{F}\) 可以是有限的或非整环的。

换言之，恒等式 \(f((x_{i})) = 0\) （当 \(x_{i}\) 遍历一个以 A 为子环的无限整环时，且可带有如下限制： \(g_{k}((x_{i})) \neq 0\) （对 \(1 \leqslant k \leqslant m\) ，其中 \(g_{k}\) 是非零多项式））的证明，蕴含当 \(x_{i}\) 遍历任意含单位元的结合交换 A-代数时的同一恒等式。

特别地，设 \(f \in \mathbf{Z}[(\mathbf{X}_i)_{i \in I}]\) 。若 \(f((x_i)) = 0\) （当 \(x_i\) 遍历 \(\mathbf{Z}\) 时，可带有限制： \(g_k((x_i)) \neq 0\) （对 \(1 \leqslant k \leqslant m\) ，其中 \(g_k\) 是 \(\mathbf{Z}[(\mathbf{X}_i)]\) 的非零元素）），则当 \(x_i\) 遍历任意交换环时，同一恒等式成立。

